% !TeX root = ../papers/nc-p3.tex
\section{特異四次核の境界行列と節点二捩れ}
\label{sec:singular-boundary-bezout}
\begin{theorem}\label{thm:singular-boundary-bezout}
四辺有限被約の標準有限Frobenius窓には、整な反対八次曲面で余因子の包が八次元のものはない。
境界Hankel元の場合、被約性を外しても点スキームのうち以下の導手直線の外にある部分の長さは16以下である。
\end{theorem}
\begin{proof}[核列の座標化]
残るのは節点・尖点での境界と節点の $t=-1$ である。
境界では四行Hankel行列の随伴を $d^2$ で割った対称外積の列を
$a(z)=a_0+a_1z+a_2z^2+a_3z^3$ とする。$p$ を特異点として
$f=N(p,p)$ は節点で $a_0a_3$、尖点で $a_0^2$ の定数倍であり、非零である。
一般の $x$ で八切断上の余因子汎関数は、$p$ と正規化上の $a$ の三根を合わせた
長さ4の部分スキームでの評価を経由する。実際 $p$ で消える $L^2$ の切断は
節点で $zS_6$、尖点で $z^2S_6$ で、そのうち $aS_3$ に上記の因子を掛けたものを
Hankel核が消す。$S_j$ は次数 $j$ の二項式空間を表し、アフィン表示でも無限遠の次数を数える。
節点二捩れでは $a$ は $P=L\eta$ の四次切断で $a_4=-a_0$、$f$ は $a_0^2$ に比例し、
$aH^0(P)$ の消滅から余因子汎関数は $a$ の四根を経由する。

いずれも $a$ は四つの独立な入力座標を与える。実際 $[a]$ を固定すると一般の可逆な $L$ 対合の
逆の線型包は高々四次元である。固定長さ4の環への $L$ の評価が同型であり、他の可逆対合は
一つの可逆対合と環の単元の乗法の合成だからである。階数2以下の $a$ なら一般繊維に平面が入り、
全二次式 $N(x)$ のその平面上の像が六次元を張ることに反する。階数3とし入力核点を $o$ とする。
正規化の一般の二点 $z,w$ の両方で消える切断がこの射影二次元線形系に存在する。
一般の対は一般切断を与える。一般切断は単純根を持ち（固定二重点は二つの独立条件であり系全体は入らない）、
二根を取る入射の切断側への繊維は有限である。固定単根があっても動く根は二つ以上ある。
この切断の入力直線上では $N(x)(z,w)=0$ である。実際一般点で余因子対合は
四評価の重み付き和で、逆二変数形式はこれらの点の基底で対角になる。
直線上の恒等式を $o$ に制限すると $N(o)(z,w)=0$ となる。一般対での消滅と
$E$ の非退化性から $N(o)=0$ となり、全二次式の無基点性に反する。
\end{proof}

\begin{proof}[境界の全ての二次式表示]
$F(z)=N(p,z)$ は一般の $a$ の三根で消え、普遍三次式 $a(z)$ で割れる。
商は入力に一次、$z$ に高々一次である。$F$ 自身の節点貼合と尖点貼合から順に
$F=a(z)(a_3+a_0 z)$、$F=a(z)(a_0-a_1z)$ を得る（全体倍率を規格化する）。
実際節点の商 $c+dz$ は $a_0c=a_3d$、尖点では $a_1c+a_0d=0$ を満たす。
次の特解をそれぞれ取る。
\begin{align*}
 N_{\rm n}^*(z,w)&=(z+w)a(z)a(w)
 +\frac{a_3(za(z)-wa(w))+a_0zw(z^4a(w)-w^4a(z))}{z-w},\\
 N_{\rm c}^*(z,w)&=a(z)a(w)-a_1(za(z)+wa(w))-a_1zw\frac{a(z)-a(w)}{z-w}.
\end{align*}
両式は対称で、定数項と最高次項、または一次項を読むと両変数で貼合を満たす。
$p$ の行が $F$ となり、相異なる二根で零となる。
$g(z)$ を節点で $z$（同次形は両端で消える二次式）、尖点で $z^2$ とする。
$N-N^*$ を $g(z)g(w)$ で割ると、各変数二次で $a$ の三根において対角な形式となる。
二項Bezout同型
\[
 B(c,d)=\frac{c(z)d(w)-c(w)d(z)}{z-w},\qquad
 \bigwedge^2 S_3\simeq\operatorname{Sym}^2 S_2
\]
の下でこの対角三空間は $a\wedge S_3$ である。実際これらの形式は対角で、核は $k a$ である。
従って入力に二次の交代テンソル $\Omega$ は $a\wedge\Omega=0$ を満たす。
座標四つのKoszul完全性で $\Omega=a\wedge Ta$ と書け、境界の全表示を得る。
\[
                  N=N^*+g(z)g(w)B(a,Ta),\qquad T\in\operatorname{End}(S_3)/kI .
\]
\end{proof}

\begin{proof}[境界の導手直線外の残余長]
階数1で $f=0$ なら $F=0$ なので、その台は節点で $C=(a_0,a_3)$、尖点で
$C=(a_0,a_1)$ という入力直線に限られる。$f\ne0$ では $F(z)F(w)/f$ を引いて
$g(z)g(w)$ で割ると $B(a,Ta-U_u(a))$ となる。ここで
\[
\begin{array}{lll}
\text{節点}& u=a_3/a_0,& U_u(a)=u(a-a_0)/z-z(a-a_3z^3)/u,\\
\text{尖点}& u=a_1/a_0,& U_u(a)=(a_1+(-u+u^2z)a)/z^2 .
\end{array}
\]
実際両辺は三根の基底で対角である。単純根 $r$ での特解の寄与は分母 $g(r)^2$ を除いた後、
節点で $a'(r)(a_3/r-a_0r^4)$、尖点で $-a'(r)a_1/r^2$ である。
これで一般の $a$ で一致し、分母を可逆にした環で恒等式となる。
Bezout同型から階数1の式は $Ta-U_u(a)=\lambda a$ とスキームとして同値である。

$u$ を完備化し、節点では $a_3=u a_0$、尖点では $a_1=u a_0$ の超平面束を
$A\subset S_3\otimes\mathcal O_{\mathbb P^1}$ とする。階数3、次数 $-1$ である。
節点では $D'=0+\infty$、尖点では $D'=2\infty$ を極因子に許せば $U:A\to S_3(D')$ は正則である。
節点の0での先頭項は $a\mapsto-za$、両型の無限遠での先頭項は $a\mapsto a/z$ である。
いずれも包含との鉛筆に固有ベクトルを持たない。直線の規約で
$\Sigma=\mathbb P(\mathcal O\oplus\mathcal O(D'))$ を取ると上の条件は
$A\otimes\mathcal O_\Sigma(-1)\to S_3(D')$ の退化となる。
$\lambda=\infty$ では包含であり、極上も退化しない。従って退化は有限である。
実際射影曲線なら $u$ の極を避けて垂直となり、その繊維でも $\lambda$ の無限遠を避けられない。
$h=c_1(\mathcal O_\Sigma(1)),\ e=[\text{繊維}]$ とすると $h^2=-2he,\int he=1$。
源の第一・第二類は $-e-3h,2eh+3h^2$、的の第一類は $8e$ であり、
有限な最大行列式の退化長は $c_2(\mathrm{target}-\mathrm{source})=16$ である。

非被約な点窓にも同じ上界を使える。退化上の直線を印す入射に対応する線型方程式の
次数2部分のKoszul複体を取る。階数が1以下落ちる一般点を除き得ることは
任意の $T$ を動かした評価が $\operatorname{Hom}(A_u,S_3)$ に全射であることから従う。
一般の $T$ では退化は横断的な16点、各余核は長さ1である。任意の $T$ でも複体の
正次数ホモロジーは有限集合上に支えられ、二次部分は長さ2の局所自由複体なので完全である。
実際左端は一般単射、次の余核の深さは滑らかな曲面上で1以上なので有限長部分を持たない。
Euler数は一定ゆえ終端余核の長さは常に16である。各ファイバーの核は高々二次元である。
三次元なら超平面 $A_u$ 全体で残余方程式が成立し、$a_0\ne0$ の稠密開上で
$N$ の階数が1となり整な八次曲面に平面が含まれるからである。導手外で $u,\lambda$ は
入力から正則に一意復元され、入射の各点窓ファイバーは入力点または直線切断の閉部分である。
直線切断長は3以下なので相対二次式がその全切断を生成し、中山の補題で長さは16以下となる。
被約窓では導手直線上も高々3点なので二十点に届かない。
\end{proof}

\begin{proof}[節点の二捩れの微分表示]
$t=-1+s$ と置き $s^2=0$ の上で二つの貼合基底を取る。逆元は $-1-s$ なので
$C=C_0+sC_1$ の一次項は交代である。滑らかな二捩れの証明と同様に
$K=\operatorname{adj}(C)/d^2=aa^{\mathsf t}+sK_1$ が存在する。
$C_0 K_1=-C_1 aa^{\mathsf t}$ から $a\wedge K_1=0$、従って
$K_1=ba^{\mathsf t}-ab^{\mathsf t}$ と線型列で書ける。
二つの基底の符号を対応させると正規化上の切断は $a\pm s b$ で、
$a_4=-a_0,\ b_4=-b_0+a_0$ である。$b$ は $a$ の固定線型写像である。
非分岐節点の二変数核に $t=-1+s$ とこれらの切断を入れた一次係数は
\[
 M(a)(z,w)=a(z)a(w)+2(z+w)\frac{b(z)a(w)-a(z)b(w)}{z-w} .
\]
これが $N$ を与えることを確かめる。一般の $a$ の単純四根を $r_i$ として、
余因子汎関数を $\sum c_i\operatorname{ev}_{r_i}$ と書く。全 $c_i\ne0$ である。
核の一次持上げの式を四根に制限すると $\sum c_i b(r_i)h(r_i)=0$ が全 $h\in H^0(P)$ に成立する。
四評価にはただ一つの関係があり、その全係数は非零である。実際任意の三根の評価は
次数1の残余線束によって独立である。また $b$ が四根全てで消えれば $a$ に比例し、
貼合 $b_4+b_0=a_0\ne0$ に反する。微分 $(h/a)\,dz/z$ の留数和は両端で相殺されるので
$c_i$ は $1/(r_i a'(r_i)b(r_i))$ に共通定数倍で等しい。
$M$ はこれらの根で対角で、対角値は $-4r_i a'(r_i)b(r_i)$ である。
よって $N$ と射影的に一致し、全二次式の原始性から定数倍である。
\end{proof}

\begin{proof}[節点二捩れの全ての被約台点]
$a,b$ から共通因子を取り除いた残余を $A,B$ とする。比例しないとして残余次数を $m>0$ とする。
$A$ が0で消えると残余核の対角 $(0,0)$ は零だが0の行は $-2B(0)A(z)$ で非零となる。
階数1に反する。無限遠でも反転で交代項の符号が変わるだけで同様である。
そこで $R=k[z]/A$ で最初の変数を評価すると、単元 $B(z)$ と全体倍率を除き
$A(w)+2z A(w)/(w-z)$ の係数が得られる。$w^m$ の係数と三角な商の係数は全 $R$ を張り、
階数1ゆえ $m=1$ である。$A=z-u$, $B=\lambda(z-u)+v$ と書けば
残余行列式は $-4v(u+v)$ なので、比例しないとき $v=-u$ である。
従って台点は次の入射で尽くされる。
\[
       a=(rz-s)g,\qquad b=\lambda(rz-s)g-sg,\qquad r g_3=s g_0 .
\]
比例する場合（$b=0$ も含む）は貼合により $a_0=a_4=0$ なので $s=0$ で覆える。
$g$ は三次式で、$b+sg$ と $a$ は同じ貼合空間に入る。
$[r:s]$ 上の $a$ の束 $A$ は次数 $-4$ である（四次切断の評価核）。
スカラーが無限遠の写像は包含なので単射であり、水平退化は一定スカラーに限られる。
この各スカラーの像は $b_0=(\lambda+1)a_0$ の平面に入る。
$a_0,b_0$ が従属なら $a_0=0$ の平面上で $a,b$ がともに $0$ で消え、
$M$ の $p$ 行が零で八次行列式がこの平面を含む。従って各平面は真である。

定理\ref{thm:singular-residual-pencils}の定数固有直線を取り除く議論がこの鉛筆にも適用できる。
各商で無限遠の包含は単射であり、各除去スカラー以外では商から元の核を一意に戻せる。
有限残余の長さは除去回数 $j=0,1,2$ に応じて $16,12-b',8-b'$ で、
$b'$ は除いた部分束の次数の負である。異なる二スカラーで $j=2$ のとき $b'\ge3$ も成立する。
両固有線が次数1なら各 $g$ は $r,s$ に依らず $g_0=g_3=0$ となり、二固有線はいずれも
$a_0=b_0=0$ の直線を与える。その線型系 $\langle g,z g\rangle$ の最大公約因子で
$g$ は定数倍まで同じになり、$r$ の係数の比較で $b(zg)=\lambda zg$ となって矛盾する。
従って同定または除去された諸平面の上の台点も数えて高々19点である
（$j=3$ なら三平面の和に入り三次単射性に反する）。前節と滑らかな場合の排除を合わせて結論を得る。
\end{proof}

\path{computations/verify_singular_boundary_bezout.py} は両境界の貼合と固有ベクトル条件を
有理入力で検算し、各型の固定標本について二次成分の階数10、良い素数で四次小行列式の階数19を確認する。
